\documentclass[11pt,letterpaper]{article}
\usepackage{amsmath,amsthm}
\usepackage{unicode-math}
\setmainfont{LibertinusSerif-Regular.otf}[
  BoldFont=LibertinusSerif-Bold.otf,
  ItalicFont=LibertinusSerif-Italic.otf,
  BoldItalicFont=LibertinusSerif-BoldItalic.otf
]
\setmathfont{LibertinusMath-Regular.otf}
\usepackage{booktabs}
\usepackage[hidelinks]{hyperref}

\setlength{\textwidth}{6.30in}
\setlength{\textheight}{9.12in}
\setlength{\headheight}{12pt}
\setlength{\headsep}{17pt}
\setlength{\footskip}{25pt}
\setlength{\oddsidemargin}{\dimexpr(\paperwidth-\textwidth)/2-1in\relax}
\setlength{\evensidemargin}{\oddsidemargin}
\setlength{\topmargin}{\dimexpr 0.96in-1in-\headheight-\headsep\relax}
\setlength{\parindent}{1.1em}
\setlength{\parskip}{2pt plus 0.5pt}
\linespread{1.065}
\raggedbottom
\clubpenalty=10000
\widowpenalty=10000
\displaywidowpenalty=10000
\setlength{\jot}{5pt}
\AtBeginDocument{%
  \setlength{\abovedisplayskip}{11pt plus 2pt minus 2pt}%
  \setlength{\belowdisplayskip}{11pt plus 2pt minus 2pt}%
  \setlength{\abovedisplayshortskip}{6pt plus 2pt minus 1pt}%
  \setlength{\belowdisplayshortskip}{9pt plus 2pt minus 1pt}%
}

\makeatletter
\def\ps@preprint{%
  \let\@mkboth\@gobbletwo
  \def\@oddhead{\normalfont\small Parking-function polytopes\hfil\thepage}%
  \let\@evenhead\@oddhead
  \let\@oddfoot\@empty
  \let\@evenfoot\@empty
}
\pagestyle{preprint}
\renewcommand{\@seccntformat}[1]{\csname the#1\endcsname.\enspace}
\renewcommand{\section}{\@startsection{section}{1}{\z@}%
  {-22pt plus -4pt minus -2pt}{8pt plus 2pt}%
  {\normalfont\fontsize{14}{17}\selectfont\bfseries\raggedright}}
\renewcommand{\subsection}{\@startsection{subsection}{2}{\z@}%
  {-15pt plus -3pt minus -2pt}{5pt plus 1pt}%
  {\normalfont\bfseries\raggedright}}
\renewenvironment{proof}[1][\proofname]{%
  \par\pushQED{\qed}\normalfont
  \topsep9pt plus 2pt minus 1pt\relax
  \trivlist\item[\hskip\labelsep\itshape #1\@addpunct{.}]\ignorespaces
}{\popQED\endtrivlist\@endpefalse}

\def\@maketitle{%
  \begingroup
  \raggedright
  {\fontsize{18}{22}\selectfont\bfseries\@title\par}%
  \vskip9pt
  {\normalsize\@author
    \thanks{AI-authored. See Section~\ref{sec:attribution} for contributions and verification.}\par}%
  \vskip3pt
  {\small\@date\par}%
  \endgroup
  \vskip17pt
}
\renewenvironment{abstract}{%
  \par\noindent\textbf{Abstract}\par\nobreak\vskip4pt
  \noindent\ignorespaces
}{\par}
\setlength{\abovecaptionskip}{0pt}
\setlength{\belowcaptionskip}{8pt}
\long\def\@makecaption#1#2{%
  \vskip\abovecaptionskip
  {\small\noindent\textbf{#1.}\enspace #2\par}%
  \vskip\belowcaptionskip
}
\makeatother

\hypersetup{
  pdftitle={Counterexamples to magic positivity and Ehrhart positivity of generalized parking-function polytopes},
  pdfauthor={Edwin Chan},
  pdfsubject={Ehrhart and magic positivity of generalized parking-function polytopes},
  pdfkeywords={Ehrhart polynomial, parking-function polytope, magic positivity}
}

\newtheoremstyle{result}
  {11pt plus 2pt minus 1pt}{9pt plus 2pt minus 1pt}
  {\itshape}{}{\bfseries}{.}{0.55em}{}
\theoremstyle{result}
\newtheorem{theorem}{Theorem}[section]
\newtheorem{lemma}[theorem]{Lemma}
\newtheorem{proposition}[theorem]{Proposition}
\newtheorem{corollary}[theorem]{Corollary}
\theoremstyle{remark}
\newtheorem{remark}[theorem]{Remark}

\newcommand{\NN}{\mathbb{N}}
\newcommand{\ZZ}{\mathbb{Z}}
\newcommand{\QQ}{\mathbb{Q}}
\newcommand{\RR}{\mathbb{R}}
\newcommand{\X}{\mathfrak{X}}
\newcommand{\vb}{\mathbf{b}}
\newcommand{\vx}{\mathbf{x}}
\newcommand{\ehr}{\operatorname{ehr}}

\begin{document}

\title{Counterexamples to magic positivity and Ehrhart positivity of generalized parking-function polytopes}
\author{Edwin Chan}
\date{November 26, 2026}

\maketitle

\begin{abstract}
We disprove Conjecture~8.1 of Hill, Luo, Trinh, and Vindas-Mel\'endez by giving a generalized parking-function polytope in dimension $14$ with a negative second magic coefficient. A second example in the same dimension has a negative second Ehrhart coefficient. The construction uses a Minkowski sum of two uniform-matroid independence polytopes. We determine the minimum possible largest entry for each sign failure in a specified two-parameter family, give a binary example in dimension $20$ exhibiting both failures, and exhibit negative mixed multivariate coefficients in dimensions $13$ and $14$.
\end{abstract}

\smallskip
{\small\noindent\textit{2020 Mathematics Subject Classification.}\enspace 52B20, 05A15, 05B35.
\par\noindent\textit{Key words and phrases.}\enspace Generalized parking-function polytope, Ehrhart polynomial, magic positivity, uniform matroid.\par}

\section{Introduction}
\label{sec:intro}

For a positive integer vector $\vb=(b_1,\dots,b_n)\in\ZZ_{>0}^n$, let $S_j:=\sum_{r=1}^j b_r$ ($1\le j\le n$) and $\Sigma_k:=\sum_{j=n-k+1}^n S_j$ ($1\le k\le n$). The \emph{generalized parking-function polytope} (or \emph{$\vb$-parking-function polytope}) $\X_n(\vb)\subset\RR^n$, introduced in~\cite{BBCDDDLMMNV,HLV24}, is the convex hull of all $\vb$-parking functions of length $n$. By~\cite[Theorem~2.3(c)]{BBCDDDLMMNV} (restated as~\cite[Theorem~2.2]{HLTV26}), it admits the inequality description
\begin{equation}\label{eq:ineq}
\X_n(\vb) =
\left\{\mathbf{z}\in\RR^n :
  \begin{aligned}
    z_i &\ge 1 && (i\in[n]),\\
    \sum_{i\in I} z_i &\le \Sigma_{|I|} && (\emptyset\ne I\subseteq[n])
  \end{aligned}
\right\},
\end{equation}
where $[n]:=\{1,\dots,n\}$ and $\NN:=\ZZ_{\ge 0}$.

For an $n$-dimensional lattice polytope $P\subset\RR^n$, write its Ehrhart polynomial in both the standard monomial basis and the magic basis $\{t^i(t+1)^{n-i}\}_{i=0}^n$:
\[
\ehr_P(t) = \sum_{i=0}^n e_i(P)\,t^i = \sum_{i=0}^n \mu_i(P)\,t^i(t+1)^{n-i}.
\]
The polytope $P$ is \emph{Ehrhart positive} if $e_i(P)\ge 0$ for all $0\le i\le n$, and \emph{magic positive}~\cite{FH24} if $\mu_i(P)\ge 0$ for all $0\le i\le n$. Magic positivity is strictly stronger than Ehrhart positivity and, by Br\"and\'en~\cite[Theorem~4.2]{Bra06}, implies that the Ehrhart $h^*$-polynomial of $P$ is real-rooted. Since $e_0(P)=\mu_0(P)=1$, comparing coefficients of $t$ and $t^2$ yields
\begin{equation}\label{eq:mu2}
\mu_1(P) = e_1(P) - n,
\qquad
\mu_2(P) = e_2(P) - (n-1)\,e_1(P) + \binom{n}{2}.
\end{equation}

Ferroni and Higashitani~\cite[Problem~4.23]{FH24} asked to what extent magic positivity holds for generalized parking-function polytopes. Hill, Luo, Trinh, and Vindas-Mel\'endez~\cite[Theorem~7.2]{HLTV26} settled the two-parameter subfamily $\vb=(a,b,\dots,b)$, proving that $\X_n(a,b,\dots,b)$ is magic positive if and only if $(n,a,b)\ne(2,1,1)$ (extending the stable-range partial permutohedron classification of Liu and Zhang~\cite[Theorem~1.3]{LZ26}), and formulated the following general conjecture.

\begin{quote}
\textbf{Conjecture~8.1 of~\cite{HLTV26}.} \emph{For every $n\ge 3$ and every $\vb\in\ZZ_{>0}^n$, the polytope $\X_n(\vb)$ is magic positive. Equivalently, $PF_2=\X_2(1,1)$ is the only $\vb$-parking-function polytope that is not magic positive.}
\end{quote}

In this paper we refute Conjecture~8.1 and show that even Ehrhart positivity fails within the family $\X_n(\vb)$.

\begin{theorem}[Refutation of Conjecture~8.1 and Ehrhart positivity]\label{thm:main}
In dimension $n=14$, consider the one-parameter family
\[
\vb(m) := (1^5,1+m,1^3,1+m,1^4) \in \ZZ_{>0}^{14}
\qquad (m\in\NN),
\]
where $1^k$ denotes $k$ consecutive entries equal to $1$.

Then the second magic coefficient of $\ehr_{\X_{14}(\vb(m))}(t)$ is the exact quadratic polynomial
\begin{equation}\label{eq:mu2-quad}
\mu_2\bigl(\X_{14}(\vb(m))\bigr) = -\frac{36952721}{3880800}\,m^2 + \frac{149416051}{55440}\,m + \frac{29211}{4}.
\end{equation}
In particular:
\begin{enumerate}
\item[(a)] At $m=286$, with $\vb=(1^5,287,1^3,287,1^4)$,
\[
\mu_2\bigl(\X_{14}(\vb)\bigr) = -\frac{8337823}{11025} < 0,
\]
whereas $\mu_2(\X_{14}(\vb(285)))=\frac{102746219}{51744}>0$. Thus $\X_{14}(1^5,287,1^3,287,1^4)$ is not magic positive, refuting Conjecture~8.1 of~\cite{HLTV26}.
\item[(b)] At $m=317$, with $\vb=(1^5,318,1^3,318,1^4)$,
\[
e_2\bigl(\X_{14}(\vb)\bigr) = -\frac{777669659}{3880800} < 0,
\qquad
\mu_2\bigl(\X_{14}(\vb)\bigr) = -\frac{11195736263}{117600} < 0,
\]
so $\X_{14}(1^5,318,1^3,318,1^4)$ is not even Ehrhart positive.
\end{enumerate}
\end{theorem}

Section~\ref{sec:decomp} gives the Minkowski decomposition of $\X_n(\vb)-\mathbf{1}$ and extracts $e_1$ and $e_2$ from the formula of Hill, Luo, Trinh, and Vindas-Mel\'endez. Section~\ref{sec:proof} derives the dimension-$14$ coefficient formulas. Section~\ref{sec:sharper} gives smaller-entry examples and establishes the two thresholds within the specified family.

\section{Minkowski decomposition and Ehrhart coefficients}
\label{sec:decomp}

For each $r\in\{1,\dots,n\}$, let $U_r:=\{z\in[0,1]^n : \sum_{i=1}^n z_i\le r\}$ be the independence polytope of the uniform matroid of rank $r$ on $[n]$, and let $P_0:=\X_n(1,\dots,1)-\mathbf{1}$ be the shifted classical parking-function polytope.

\begin{proposition}[Minkowski decomposition]\label{prop:mink}
For every $\vb\in\ZZ_{>0}^n$, let $x_i:=b_i-1\in\NN$ ($1\le i\le n$). Then
\begin{equation}\label{eq:mink}
P(\vx) := \X_n(\vb) - \mathbf{1} = P_0 + \sum_{r=1}^n x_r\, U_{n-r+1}.
\end{equation}
\end{proposition}

\begin{proof}
By~\eqref{eq:ineq}, $P(\vx)$ is the set of points $z\in\RR_{\ge 0}^n$ satisfying $\sum_{i\in I} z_i \le \Sigma_{|I|}-|I|$ for all $\emptyset\ne I\subseteq[n]$. Substituting $S_j=j+\sum_{r=1}^j x_r$ into $\Sigma_k-k$ ($1\le k\le n$) gives
\[
\Sigma_k - k = \left(\sum_{j=n-k+1}^n j - k\right) + \sum_{j=n-k+1}^n \sum_{r=1}^j x_r.
\]
Each $x_r$ appears in the inner sum precisely for $j\in\{\max(r,n-k+1),\dots,n\}$, a total of $n-\max(r,n-k+1)+1=\min(n-r+1,k)$ times. Setting $f_0(I):=\sum_{j=n-|I|+1}^n j - |I|$ and $f_r(I):=\min(n-r+1,|I|)$, we therefore have
\[
\Sigma_{|I|}-|I| = f_0(I) + \sum_{r=1}^n x_r\, f_r(I)
\qquad\text{for all } I\subseteq[n].
\]
Writing $f_0(k)$ for its value on $k$-element subsets, we have $f_0(k)-f_0(k-1)=n-k\ge 0$ is nonincreasing in $k$, so $f_0$ is a monotone submodular rank function cutting out $P_0$; likewise each $f_r$ is the rank function of the uniform matroid $U_{n-r+1}$. All these functions are normalized, nondecreasing, and submodular on the same ground set. Schrijver's Minkowski sum identity~\cite[p.~781, Theorem~44.6]{Sch03} gives $P_{f+g}=P_f+P_g$ for their polymatroids, and $P_{\lambda f}=\lambda P_f$ for $\lambda\ge 0$. Iterating the identity yields~\eqref{eq:mink}.
\end{proof}

Next, recall from~\cite[Section~5 and Theorem~5.9]{HLTV26} the signed simplex coefficients
\begin{equation}\label{eq:wk}
w_1(\vb) := b_1 - 1,
\qquad
w_k(\vb) := \sum_{j=0}^{k-2}(-1)^j\binom{k-2}{j}\, b_{k-j} \quad (2\le k\le n),
\end{equation}
and the draconian Ehrhart formula, obtained there from Postnikov's lattice-point formula~\cite[Theorem~11.3]{Pos09} by extending to the signed coefficients:
\begin{equation}\label{eq:thm59}
\ehr_{\X_n(\vb)}(t) = \sum_{\alpha\in\mathcal{D}(n)} \prod_{\emptyset\ne A\subseteq[n]} \binom{w_{|A|}(\vb)\,t + \alpha_A - 1}{\alpha_A},
\end{equation}
where $\mathcal{D}(n)$ consists of all sequences $\alpha=(\alpha_A)_{\emptyset\ne A\subseteq[n]}\in\NN^{2^n-1}$ satisfying the Hall condition $\sum_{A\in E}\alpha_A\le \bigl|\bigcup_{A\in E} A\bigr|$ for every nonempty family $E$ of nonempty subsets of $[n]$. Denote by $H_m:=\sum_{j=1}^m\frac{1}{j}$ the $m$-th harmonic number (with $H_0:=0$).

\begin{samepage}
\begin{lemma}[Closed forms for $e_1$ and $e_2$]\label{lem:e1e2}
For every $\vb\in\ZZ_{>0}^n$, with $w_k:=w_k(\vb)$ defined by~\eqref{eq:wk},
\begin{equation}\label{eq:e1-closed}
e_1\bigl(\X_n(\vb)\bigr) = \sum_{k=1}^n \binom{n}{k} H_k\, w_k,
\end{equation}
\begin{equation}\label{eq:e2-closed}
\begin{aligned}
e_2\bigl(\X_n(\vb)\bigr) &= \sum_{k=1}^n \binom{n}{k} w_k^2 \sum_{a=1}^k \frac{H_{a-1}}{a} \\
  &\quad + \frac{1}{2}\sum_{(k,\ell,u)\in\mathcal{T}_n}
    \frac{n!\,w_k w_\ell}{(n-u)!\,(u-\ell)!\,(u-k)!\,(k+\ell-u)!}
    \sum_{a=1}^{\min(k,u-1)}\frac{H_{\min(\ell,u-a)}}{a},
\end{aligned}
\end{equation}
where $\mathcal{T}_n:=\{(k,\ell,u)\in[n]^3 : \max(k,\ell)\le u\le\min(n,k+\ell),\; (k,\ell,u)\neq(u,u,u)\}$.
\end{lemma}
\end{samepage}

\begin{proof}
For every integer $a\ge 1$, expanding $\binom{y+a-1}{a}=\frac{y(y+1)\cdots(y+a-1)}{a!}$ around $y=0$ gives
\[
\binom{y+a-1}{a} = \frac{1}{a}\,y + \frac{H_{a-1}}{a}\,y^2 + O(y^3),
\]
while $\binom{y-1}{0}=1$. Substituting $y=w_{|A|}t$ into~\eqref{eq:thm59}:
\begin{itemize}
\item Only sequences $\alpha\in\mathcal{D}(n)$ supported on a single subset $A$ ($|A|=k$, $1\le \alpha_A\le k$) contribute to $[t^1]$, yielding $\sum_{k=1}^n\binom{n}{k}\sum_{a=1}^k\frac{w_k}{a}=\sum_{k=1}^n\binom{n}{k}H_k w_k$, which is~\eqref{eq:e1-closed}.
\item Sequences $\alpha\in\mathcal{D}(n)$ contributing to $[t^2]$ have support size either $1$ or $2$. Support size $1$ ($|A|=k$, $1\le a=\alpha_A\le k$) contributes $\frac{H_{a-1}}{a}w_k^2$ per subset $A$, giving the first sum in~\eqref{eq:e2-closed}. Support size $2$ consists of unordered distinct pairs $\{A,B\}$ of nonempty subsets with $|A|=k$, $|B|=\ell$, and $|A\cup B|=u$ (excluding $k=\ell=u$, which forces $A=B$). For such a pair, the Hall conditions reduce to $1\le a\le k$, $1\le b\le\ell$, and $a+b\le u$, i.e., $1\le a\le\min(k,u-1)$ and $1\le b\le\min(\ell,u-a)$, contributing
\[
\begin{aligned}
\sum_{a=1}^{\min(k,u-1)}\sum_{b=1}^{\min(\ell,u-a)}\frac{w_k w_\ell}{ab}
  &= w_k w_\ell \sum_{a=1}^{\min(k,u-1)}\frac{H_{\min(\ell,u-a)}}{a}.
\end{aligned}
\]
The number of ordered pairs $(A,B)$ of nonempty subsets of $[n]$ with $(|A|,|B|,|A\cup B|)=(k,\ell,u)$ is
\[
\binom{n}{u}\frac{u!}{(u-\ell)!(u-k)!(k+\ell-u)!}
=\frac{n!}{(n-u)!(u-\ell)!(u-k)!(k+\ell-u)!}.
\]
Dividing by $2$ to count unordered pairs gives the second sum in~\eqref{eq:e2-closed}. \qedhere
\end{itemize}
\end{proof}

\section{Coefficient formulas in dimension 14}
\label{sec:proof}

Set $\vb(a,b)=(1^5,1+a,1^3,1+b,1^4)$ for $a,b\in\NN$ and write $e_i(a,b)$ and $\mu_i(a,b)$ for the coefficients of $\X_{14}(\vb(a,b))$. In~\eqref{eq:wk}, the all-ones vector contributes $w_2=1$ and zero at every other index. The shift at position $6$ contributes $a(-1)^{k-6}\binom{k-2}{4}$ for $k\ge 6$, and the shift at position $10$ contributes $b(-1)^{k-10}\binom{k-2}{8}$ for $k\ge 10$. The shifts contribute zero below their respective positions. In particular,~\eqref{eq:e1-closed} gives
\begin{equation}\label{eq:e1-bivariate}
e_1(a,b)=\frac{273}{2}+\frac{4613814a+3546355b}{360360}.
\end{equation}
The constant here is $\binom{14}{2}H_2$. The numerators of the two linear terms follow by summing $\binom{14}{k}H_k$ times the indicated signed binomial weights.

For $D=50450400$, substitution in~\eqref{eq:e2-closed} and then~\eqref{eq:mu2} gives
\begin{equation}\label{eq:two-quadratics}
\begin{aligned}
q(a,b) &= 4092523257a^2-7085907205ab+2512998575b^2,\\
D e_2(a,b) &= q(a,b)+84978997740a+65841116250b+453359907000,\\
D\mu_2(a,b) &= q(a,b)+76581856260a+59386750150b+368426658600.
\end{aligned}
\end{equation}
For example, the mixed coefficient of $e_2$ is the difference
$e_2(1,1)-e_2(1,0)-e_2(0,1)+e_2(0,0)$ from~\eqref{eq:e2-closed}.
The coefficients of $a^2$ and $b^2$ are half the second differences on the respective axes, and the linear terms follow after subtracting the constant and quadratic terms. This recovers the quadratics because~\eqref{eq:wk} is affine in $a,b$ and~\eqref{eq:e2-closed} has total degree at most two. The companion verifier also derives their coefficients by counting lattice points at the six sample pairs in this calculation and interpolating.

\begin{proof}[Proof of Theorem~\ref{thm:main}]
Setting $a=b=m$ in~\eqref{eq:two-quadratics} gives~\eqref{eq:mu2-quad}, which is positive at $m=285$ and negative at $m=286$ with the stated values. From~\eqref{eq:e1-bivariate},
\begin{equation}\label{eq:e1-diagonal}
e_1(m,m)=\frac{273}{2}+\frac{8160169}{360360}\,m.
\end{equation}
At $m=317$, \eqref{eq:mu2-quad} gives $\mu_2(317,317)=-11195736263/117600$. Consequently,
\[
\begin{aligned}
e_2(317,317) &= \mu_2(317,317)+13e_1(317,317)-91\\
            &= -\frac{777669659}{3880800}<0.
\end{aligned}
\qedhere
\]
\end{proof}

The leading coefficient of~\eqref{eq:mu2-quad} has a geometric interpretation. The Bernstein--McMullen polynomiality theorem~\cite{McM78} gives a polynomial for the lattice-point count in a Minkowski sum of integer dilates of lattice polytopes. If the sum is itself dilated by $t$, its $t^j$ coefficient is the homogeneous degree-$j$ part of that polynomial, evaluated at the dilation weights. In particular,
\[
P(\mathbf{x}(m))=P_0+m(U_5+U_9),
\]
by Proposition~\ref{prop:mink}. The coefficient of $m^2$ in $e_2(m,m)$ therefore comes from $U_5+U_9$ alone. As $e_1(m,m)$ is affine, the same coefficient occurs in $\mu_2(m,m)$. Hence
\begin{equation}\label{eq:U5U9}
[t^2]\ehr_{U_5+U_9}(t)=-\frac{36952721}{3880800}<0.
\end{equation}
The verifier also interpolates this coefficient from lattice counts. For each dilate, the constraints are $z_i\ge 0$ and bounds on the sums of the $k$ largest coordinates, given by the dilation factor times $\min(k,5)+\min(k,9)$. Grouping equal coordinates gives an integer dynamic program, with binomial factors for assigning coordinate positions.

\section{Smaller entries and mixed coefficients}
\label{sec:sharper}

Table~\ref{tab:counterexamples} gives the equal-weight dimension-$14$ cases, two unequal-weight cases, and examples with entries in $\{1,3\}$ in dimension $18$ and $\{1,2\}$ in dimension $20$.

\begin{table}[!ht]
\centering
\begin{minipage}{0.90\linewidth}
\caption{Exact second Ehrhart and magic coefficients.}
\label{tab:counterexamples}
\renewcommand{\arraystretch}{1.38}
\setlength{\tabcolsep}{12pt}
\begin{tabular*}{\linewidth}{@{\extracolsep{\fill}}rlrr@{}}
\toprule
$n$ & Parameter vector $\vb$ & $e_2$ & $\mu_2$ \\
\midrule
$14$ & $(1^5,287,1^3,287,1^4)$ & $\dfrac{3753767953}{44100}$ & $-\dfrac{8337823}{11025}$ \\
\addlinespace[7pt]
$14$ & $(1^5,318,1^3,318,1^4)$ & $-\dfrac{777669659}{3880800}$ & $-\dfrac{11195736263}{117600}$ \\
\addlinespace[7pt]
$14$ & $(1^5,187,1^3,228,1^4)$ & $\dfrac{3101451036227}{50450400}$ & $-\dfrac{10491632153}{50450400}$ \\
\addlinespace[7pt]
$14$ & $(1^5,208,1^3,253,1^4)$ & $-\dfrac{5255462383}{5605600}$ & $-\dfrac{35322635893}{509600}$ \\
\addlinespace[7pt]
$18$ & $(1^6,3,1^5,3,1^5)$ & $\dfrac{22435315633}{6306300}$ & $-\dfrac{3708081001}{3153150}$ \\
\addlinespace[7pt]
$20$ & $(1^7,2,1^5,2,1^6)$ & $-\dfrac{10007742397957}{908107200}$ & $-\dfrac{260316261282989}{15437822400}$ \\
\bottomrule
\end{tabular*}
\end{minipage}
\end{table}

\subsection{Minimality within the two-parameter family}
For $\vb(a,b)$ the largest entry is $1+\max(a,b)$. Thus a largest entry at most $L+1$ is equivalent to $0\le a,b\le L$. To find the minimum of either coefficient on such a box, evaluate the corresponding numerator in~\eqref{eq:two-quadratics} at every integer pair in the box and take the least value. The denominator $D$ is positive, so the search uses integers throughout. The last boxes with a positive minimum give
\[
\min_{0\le a,b\le 226}\mu_2(a,b)=\frac{653047639}{300300},
\qquad
\min_{0\le a,b\le 251}e_2(a,b)=\frac{12800604359}{7207200},
\]
attained at $(186,226)$ and $(207,251)$, respectively.

The next boxes contain the negative values $\mu_2(186,227)$ and $e_2(207,252)$ in Table~\ref{tab:counterexamples}. The minimum largest entries for these two sign failures are therefore $228$ and $253$, respectively, within $\vb(a,b)$.

\begin{samepage}
\begin{proposition}[Multivariate magic negativity in dimensions $13$ and $14$]\label{prop:multiv}
Let
\[
M_n(\vx,y):=(1-y)^n\ehr_{P(\vx)}\bigl(\frac{y}{1-y}\bigr)\in\QQ[x_1,\dots,x_n,y]
\]
be the shifted multivariate magic polynomial of $\X_n(\vb)$ in $x_i=b_i-1$. Then
\begin{align*}
[x_6 x_{10} y^2]\,M_{13} &= [t^2]\bigl(\ehr_{U_4+U_8} - \ehr_{U_4} - \ehr_{U_8}\bigr) = -\frac{842231}{69300} < 0, \\
[x_7 x_{11} x_{12} y^3]\,M_{14} &= [abc]\,\#\bigl((a U_3 + b U_4 + c U_8)\cap\ZZ^{14}\bigr) = -\frac{125057}{450} < 0.
\end{align*}
\end{proposition}
\end{samepage}

\begin{proof}
In $M_n$, the $y^d$ coefficient is $e_d$ plus a linear combination of $e_j$ for $j<d$. Each $e_j$ has parameter degree at most $j$ by multivariate Ehrhart polynomiality. Thus a squarefree monomial of degree $d$ in $[y^d]M_n$ comes from the homogeneous degree-$d$ part of $e_d$, with the weight of $P_0$ set to zero. Evaluating the marked parameters on $\{0,1\}^d$ and taking inclusion--exclusion extracts exactly that monomial. Direct lattice counts at integer dilations of the corresponding uniform-matroid sums give the displayed values. The companion verifier also extracts them from the draconian formula.
\end{proof}

\begin{samepage}
\section{Contributions and verification}
\label{sec:attribution}

GPT-6.2 Sol (ChatGPT) found the dimension-$14$ magic-positivity counterexample and the negative mixed cubic coefficient. Gemini 4 Neon (Google Antigravity) developed the low-degree coefficient formulas, found the additional examples, and wrote the original verifier and manuscript. The human submitter selected the problem, prompted the systems, and coordinated the research.

The companion \href{https://github.com/CosmicCrusader23/mathslopslopslop/blob/284813d9a069ae81b3510376332afc7060895be4/disproved/hill-luo-trinh-vindas-melendez-conj-8.1-counterexample/verify.py}{Python verifier} asserts the exact coefficients and signs in Table~\ref{tab:counterexamples}, the mixed coefficients, and the finite search bounds. It compares draconian extraction with lattice counts, using parameter interpolation for the four large-entry dimension-$14$ rows and direct counting at the displayed dimension-$18$ and dimension-$20$ vectors. The manuscript is unrefereed and has no Lean certification.
\end{samepage}

\begin{thebibliography}{99}
\small\raggedright
\setlength{\itemsep}{3pt plus 1pt}

\bibitem{BBCDDDLMMNV}
M.~M.~Bayer, S.~Borgwardt, T.~Chambers, S.~Daugherty, A.~Dawkins, D.~Deligeorgaki, H.-C.~Liao, T.~McAllister, A.~Morrison, G.~Nelson, and A.~R.~Vindas-Mel\'endez,
\emph{Combinatorics of generalized parking-function polytopes},
\href{https://doi.org/10.1007/s00454-025-00770-1}{Discrete Comput. Geom.} \textbf{76} (2026), no.~1, 339--377.

\bibitem{Bra06}
P.~Br\"and\'en,
\emph{On linear transformations preserving the P\'olya frequency property},
\href{https://doi.org/10.1090/S0002-9947-06-03856-6}{Trans. Amer. Math. Soc.} \textbf{358} (2006), no.~8, 3697--3716.

\bibitem{FH24}
L.~Ferroni and A.~Higashitani,
\emph{Examples and counterexamples in Ehrhart theory},
\href{https://doi.org/10.4171/EMSS/86}{EMS Surv. Math. Sci.}, published online (2024).
\href{https://arxiv.org/abs/2307.10852}{arXiv:2307.10852}.

\bibitem{HLV24}
M.~Hanada, J.~Lentfer, and A.~R.~Vindas-Mel\'endez,
\emph{Generalized parking function polytopes},
\href{https://doi.org/10.1007/s00026-023-00671-1}{Ann. Comb.} \textbf{28} (2024), no.~2, 575--613.

\bibitem{HLTV26}
C.~Hill, A.~Luo, V.~Trinh, and A.~R.~Vindas-Mel\'endez,
\emph{Lattice slices, Ehrhart polynomials, and magic positivity of generalized parking-function polytopes},
\href{https://arxiv.org/abs/2607.15503v1}{arXiv:2607.15503v1}, 2026.

\bibitem{LZ26}
F.~Liu and Z.~Zhang,
\emph{Magic positivity for the Ehrhart polynomials of partial permutohedra},
\href{https://arxiv.org/abs/2607.03854v1}{arXiv:2607.03854v1}, 2026.

\bibitem{McM78}
P.~McMullen,
\emph{Lattice invariant valuations on rational polytopes},
\href{https://doi.org/10.1007/BF01226481}{Arch. Math.} \textbf{31} (1978), 509--516.

\bibitem{Pos09}
A.~Postnikov,
\emph{Permutohedra, associahedra, and beyond},
\href{https://doi.org/10.1093/imrn/rnn153}{Int. Math. Res. Not. IMRN} (2009), no.~6, 1026--1106.

\bibitem{Sch03}
A.~Schrijver,
\emph{Combinatorial optimization: polyhedra and efficiency}, Vol.~B,
Algorithms and Combinatorics \textbf{24}, Springer, Berlin, 2003.

\end{thebibliography}

\end{document}
